% !TEX program = pdflatex
\documentclass[11pt]{article}

\usepackage[T1]{fontenc}
\usepackage{lmodern}
\usepackage{microtype}
\usepackage[a4paper,margin=28mm,headheight=14pt]{geometry}
\usepackage{amsmath,amssymb,amsthm,mathtools}
\usepackage{booktabs,array}
\usepackage{xcolor}
\usepackage{xurl}
\usepackage{hyperref}
\usepackage{fancyhdr}

\definecolor{linkblue}{RGB}{24,78,119}
\hypersetup{
  colorlinks=true,
  linkcolor=linkblue,
  urlcolor=linkblue,
  pdftitle={Analytic Supplement for the Dirichlet Polya Inequality on Three-Dimensional Spherical Shells},
  pdfsubject={Ratio-sharp global defect proof}
}

\pagestyle{fancy}
\fancyhf{}
\fancyhead[L]{\small Spherical-shell P\'olya inequality}
\fancyhead[R]{\small Analytic supplement}
\fancyfoot[C]{\thepage}
\setlength{\emergencystretch}{3em}

\newtheorem{theorem}{Theorem}
\newtheorem{proposition}[theorem]{Proposition}
\newtheorem{lemma}[theorem]{Lemma}
\theoremstyle{remark}
\newtheorem{remark}[theorem]{Remark}

\title{Analytic Supplement for\\
A Dirichlet P\'olya Inequality on\\
Three-Dimensional Spherical Shells}
\author{Anonymous for peer review}
\date{17 July 2026}

\begin{document}
\maketitle

\begin{abstract}
This support volume contains the detailed global defect proof used for
$0<\rho\leq39/50$.  It proves the ratio-sharp angular estimate, the
tangent-envelope retained-radial estimate, reduction to a universal ball
quarter-layer, and the exact scalar convexity closure.  Together with the
phase bridge, no-mode comparison, and optical theorem in the main paper,
this is the complete analytic implication chain for the low- and middle-ratio
region.  Every estimate used below is proved analytically; no program,
floating-point computation, or interval certificate is a theorem premise.
\end{abstract}

\paragraph{Reading convention.}
The detailed argument is Part~I below.  In the repository it is maintained in
\path{analytic/ratio-sharp-global-closure-body.tex} and input directly; the
flattened review source inlines the same text.  No precompiled PDF is required.
A summarized form is included in the main paper.  Every displayed rational
bound below is part of the proof; numerical replay is not.

\section*{Live dependency graph}
\addcontentsline{toc}{section}{Live dependency graph}

\begin{center}
\small
\renewcommand{\arraystretch}{1.15}
\begin{tabular}{@{}p{3.3cm}p{5.1cm}p{3.2cm}@{}}
\toprule
result&premises&location\\
\midrule
phase-to-proxy bound&
separation and strict Bessel-phase enclosure&
main paper\\
no-mode region&
one-dimensional Dirichlet comparison&
main paper\\
global low/middle theorem&
exact layer cake, ratio-sharp blocks, tangent envelope,
ball quarter-layer, scalar convexity&
Part I below\\
outer optical theorem&
product and complementary-action screens&
main paper\\
shell theorem&
the preceding disjoint owners and dilation&
main paper\\
\bottomrule
\end{tabular}
\end{center}

\renewcommand{\thesection}{\Roman{section}}
\tableofcontents
\clearpage

\section{Ratio-sharp global defect closure}
\label{part:ratio}

% BEGIN mechanically inlined support source
% Shared body: included by the analytic supplement and standalone support module.

Let
\[
 A_{\rho,1}:=\{x\in\mathbb R^3:\rho<|x|<1\}.
\]
Fix $0<\rho<1$ and $K>0$.  For $a>0$, define the zero-extended ball
action
\begin{equation}
 G_a(z):=
 \begin{cases}
 \displaystyle \frac1\pi\left(\sqrt{a^2-z^2}
      -z\arccos\frac za\right),&0\leq z<a,\\[4pt]
 0,&z\geq a,
 \end{cases}
 \label{rs:eq:G}
\end{equation}
and put
\begin{equation}
 A(z):=G_K(z)-G_{\rho K}(z),
 \qquad
 T:=A(0)=\frac{(1-\rho)K}{\pi}.
 \label{rs:eq:A-T}
\end{equation}
Write
\[
 [u]_<:=\#\{n\in\mathbb N:n<u\}.
\]
The phase reduction in the main paper supplies
\begin{equation}
 N_D(A_{\rho,1},K^2)\leq P(\rho,K),
 \qquad
 P(\rho,K):=
 \sum_{\substack{\nu=\ell+1/2<K\\ \ell\in\mathbb N_0}}
 2\nu\,[A(\nu)+\tfrac14]_<.
 \label{rs:eq:phase-proxy}
\end{equation}
Indeed, the channel phase is strictly smaller than $A(\nu)+1/4$, so at an
integral proxy wall the last radial integer remains excluded.  The
continuous payment is
\begin{equation}
 W(\rho,K):=\int_0^K2zA(z)\,dz
 =\frac{2(1-\rho^3)K^3}{9\pi}.
 \label{rs:eq:W}
\end{equation}

\subsubsection*{Action geometry and the exact defect}

Differentiating \eqref{rs:eq:G} gives
\begin{equation}
 q(z):=-A'(z)=\frac1\pi
 \begin{cases}
 \displaystyle
 \arccos\frac zK-\arccos\frac z{\rho K},&0<z<\rho K,\\[6pt]
 \displaystyle \arccos\frac zK,&\rho K\leq z<K.
 \end{cases}
 \label{rs:eq:q}
\end{equation}
Thus $A$ decreases strictly from $T$ to $0$.  Let
$R:[0,T]\to[0,K]$ be its decreasing inverse and set
\begin{equation}
 F(t):=R(t)^2,
 \qquad g(t):=-F'(t),
 \qquad \tau:=A(\rho K)=KG_1(\rho).
 \label{rs:eq:inverse-data}
\end{equation}
Write throughout
\[
 \theta:=\arccos\rho.
\]

\begin{lemma}[Shape and endpoint data]
\label{rs:lem:shape}
With this notation, one has
\begin{equation}
 0\leq q(z)\leq\frac{\theta}{\pi},
 \qquad
 q\text{ increases on }(0,\rho K),
 \qquad q\text{ decreases on }(\rho K,K),
 \label{rs:eq:q-cap}
\end{equation}
and
\begin{equation}
 g\text{ decreases on }(0,\tau),
 \qquad g\text{ increases on }(\tau,T).
 \label{rs:eq:g-shape}
\end{equation}
Moreover,
\begin{equation}
 F(\tau)=\rho^2K^2,
 \qquad
 g(\tau)=\frac{2\pi\rho K}{\theta},
 \qquad
 g(T)=\frac{2\pi\rho K}{1-\rho}.
 \label{rs:eq:g-data}
\end{equation}
The two formulas for $g$ agree at $t=\tau$, and
$g(t)=O(t^{-1/3})$ as $t\downarrow0$.
\end{lemma}

\begin{proof}
Direct differentiation gives
\[
 q'(z)=\frac1\pi\left(
 \frac1{\sqrt{\rho^2K^2-z^2}}-
 \frac1{\sqrt{K^2-z^2}}\right)>0
 \quad(0<z<\rho K),
\]
whereas
\[
 q'(z)=-\frac1{\pi\sqrt{K^2-z^2}}<0
 \quad(\rho K<z<K).
\]
Both one-sided values at $z=\rho K$ equal $\theta/\pi$, proving
\eqref{rs:eq:q-cap}.

If $z=R(t)>\rho K$, then
\[
 g(t)=\frac{2\pi z}{\arccos(z/K)}.
\]
Its monotonicity follows from the explicit derivative
\[
 \frac{d}{dz}\left(\frac{z}{\arccos(z/K)}\right)
 =\frac{\arccos(z/K)+z/\sqrt{K^2-z^2}}
        {\arccos(z/K)^2}>0.
\]
Since $z$ decreases with $t$, $g$ decreases on $(0,\tau)$.  On the inner
branch put
\[
 \delta(z):=\arccos\frac zK-\arccos\frac z{\rho K}.
\]
The identity
\[
\frac{\delta(z)}z
 =\int_\rho^1\frac{ds}{s\sqrt{s^2K^2-z^2}}
\]
has derivative
\[
 \frac{d}{dz}\left(\frac{\delta(z)}z\right)
 =\int_\rho^1
 \frac{z\,ds}{s(s^2K^2-z^2)^{3/2}}>0.
\]
Hence $\delta(z)/z$ increases with $z$, and
$g(t)=2\pi z/\delta(z)$ increases as $t$ increases and $z$ decreases.
The endpoint values follow by substitution and the limit $z\downarrow0$.
Finally, writing $z=K\cos u$ on the outer branch gives
$t\asymp u^3$ and $g\asymp u^{-1}$.
\end{proof}

Put
\begin{equation}
 x_n:=n-\frac14,
 \qquad
 M(r):=\#\{\ell\in\mathbb N_0:\ell+\tfrac12<r\},
 \qquad
 N:=\#\{n\geq1:x_n<T\}.
 \label{rs:eq:grids}
\end{equation}
For an active $n$, write $r_n:=R(x_n)$ and $m_n:=M(r_n)$.

\begin{lemma}[Strict layer cake]
\label{rs:lem:layer-cake}
For every $0<\rho<1$ and $K>0$,
\begin{equation}
 P(\rho,K)=\sum_{x_n<T}m_n^2,
 \qquad
 W(\rho,K)=\int_0^T F(t)\,dt.
 \label{rs:eq:layer-cake}
\end{equation}
Both identities retain every radial and angular equality wall.
\end{lemma}

\begin{proof}
For a half-integer $\nu=\ell+1/2$,
\[
 [A(\nu)+\tfrac14]_<
 =\#\{n\geq1:x_n<A(\nu)\}.
\]
Since $A$ and $R$ are decreasing, this is equivalent to
$\nu<R(x_n)$.  Interchanging the finite sums and using
$\sum_{\ell=0}^{m-1}(2\ell+1)=m^2$ proves the first identity.  At
$R(x_n)=m+1/2$, strict angular counting gives $M(R(x_n))=m$.
The second identity is the area formula
$\int_0^TR(t)^2dt=\int_0^K2zA(z)dz$.
\end{proof}

Define
\begin{equation}
 \mathcal E_{\rm ang}:=\sum_{x_n<T}(m_n^2-r_n^2),
 \qquad
 \mathcal D_{\rm rad}:=\int_0^TF(t)\,dt-\sum_{x_n<T}F(x_n).
 \label{rs:eq:defects}
\end{equation}
Then
\begin{equation}
 W-P=\mathcal D_{\rm rad}-\mathcal E_{\rm ang}.
 \label{rs:eq:master-defect}
\end{equation}

\subsubsection*{The ratio-sharp angular payment}

\begin{proposition}[Ratio-sharp angular blocks]
\label{rs:prop:angular}
For $N\geq1$,
\begin{equation}
 \mathcal E_{\rm ang}
 <\frac{(1-\rho^2)K^2}{6}-N d(\rho),
 \qquad
 d(\rho):=\frac{3\pi}{8\arccos\rho}-\frac14.
 \label{rs:eq:angular-sharp}
\end{equation}
The inequality is strict at every common radial--angular wall.  If $N=0$,
then $\mathcal E_{\rm ang}=0$.
\end{proposition}

\begin{proof}
Fix an active $n$ and $t\in[n-1,x_n]$.  Since $R(t)\geq r_n$,
\[
 x_n-t=A(r_n)-A(R(t))
 =\int_{r_n}^{R(t)}q(s)\,ds
 \leq\frac{\theta}{\pi}(R(t)-r_n).
\]
Integration over the left block of length $3/4$ yields
\begin{equation}
 r_n\leq\frac43\int_{n-1}^{x_n}R(t)\,dt
       -\frac{3\pi}{8\theta}.
 \label{rs:eq:rn-block}
\end{equation}
Strict angular counting gives $m_n<r_n+1/2$, including when $r_n$ is a
half-integer, and hence
\begin{equation}
 m_n^2-r_n^2<r_n+\frac14
 \leq\frac43\int_{n-1}^{x_n}R(t)\,dt-d(\rho).
 \label{rs:eq:one-angular-block}
\end{equation}
The blocks are disjoint subsets of $[0,T]$.  Therefore
\[
 \mathcal E_{\rm ang}
 <\frac43\int_0^TR(t)\,dt-Nd(\rho).
\]
The inverse area identity and scaling give
\[
 \int_0^TR(t)\,dt=\int_0^KA(z)\,dz
 =\frac{(1-\rho^2)K^2}{8},
\]
which proves \eqref{rs:eq:angular-sharp}.
\end{proof}

\subsubsection*{The tangent-envelope radial payment}

Let
\begin{equation}
 C_0(t):=\#\{n\geq1:x_n<t\},
 \quad w(t):=t-C_0(t),
 \quad \mathfrak W(t):=\int_0^tw(s)\,ds,
 \quad \Psi(t):=\mathfrak W(t)-\frac t4.
 \label{rs:eq:sawtooth}
\end{equation}
On $0\leq t\leq3/4$,
\begin{equation}
 \mathfrak W(t)=\frac{t^2}{2},
 \qquad
 \Psi(t)=\frac12\left(t-\frac14\right)^2-\frac1{32}.
 \label{rs:eq:first-cell}
\end{equation}
On every later unit cell, direct integration gives
\begin{equation}
 -\frac1{32}\leq\Psi(t)\leq\frac3{32},
 \qquad
 \mathfrak W(t)\geq\frac t4-\frac1{32}.
 \label{rs:eq:sawtooth-bounds}
\end{equation}

\begin{lemma}[Retained radial remainder]
\label{rs:lem:retained}
For all $0<\rho<1$ and $K>0$,
\begin{equation}
 \mathcal D_{\rm rad}
 \geq\frac{F(\tau)}4-\frac{g(T)}8
 +g(\tau)\left(\frac\tau4+\frac3{32}\right)
 -\int_{(0,\tau)}\mathfrak W(t)\,dg(t).
 \label{rs:eq:retained-radial}
\end{equation}
\end{lemma}

\begin{proof}
Distributionally, $dw=dt-dC_0$.  Stieltjes integration by parts gives
$\mathcal D_{\rm rad}=\int_0^Tw(t)g(t)dt$; the improper boundary term
vanishes because $\mathfrak W(t)g(t)=O(t^{5/3})$ at zero.  On the
decreasing branch,
\[
 \int_0^\tau wg
 =\mathfrak W(\tau)g(\tau)
  -\int_{(0,\tau)}\mathfrak W\,dg.
\]
On the increasing branch use $w=1/4+\Psi'$,
$\int_\tau^Tg=F(\tau)$, and integrate by parts.  Since
$\mathfrak W(\tau)=\tau/4+\Psi(\tau)$, the interface value cancels and
\[
 \mathcal D_{\rm rad}
 =\frac{F(\tau)}4+\frac{\tau g(\tau)}4
  +\Psi(T)g(T)-\int_{(\tau,T)}\Psi\,dg
  -\int_{(0,\tau)}\mathfrak W\,dg.
\]
Because $dg$ is positive on $(\tau,T)$,
\[
 \Psi(T)g(T)-\int_{(\tau,T)}\Psi\,dg
 \geq-\frac{g(T)}8+\frac{3g(\tau)}{32}.
\]
This proves \eqref{rs:eq:retained-radial}.  Continuity of $g$ and
$\mathfrak W$ handles an interface on a sawtooth wall without an atom.
\end{proof}

Define the $C^1$ tangent envelope
\begin{equation}
 L_\sharp(t):=
 \begin{cases}
 t^2/2,&0\leq t\leq1/4,\\[2pt]
 t/4-1/32,&t\geq1/4,
 \end{cases}
 \qquad
 L_\sharp'(t)=\min\{t,1/4\}.
 \label{rs:eq:L-sharp}
\end{equation}
Equations \eqref{rs:eq:first-cell}--\eqref{rs:eq:sawtooth-bounds} show
\begin{equation}
 0\leq L_\sharp(t)\leq\mathfrak W(t)
 \qquad(t\geq0).
 \label{rs:eq:L-order}
\end{equation}
On $1/4\leq t\leq3/4$ the difference is
$\tfrac12(t-1/4)^2$; after $3/4$ use
\eqref{rs:eq:sawtooth-bounds}.

Assume now
\begin{equation}
 0<\rho\leq\frac{39}{50},
 \qquad K\geq\frac\pi{1-\rho}.
 \label{rs:eq:global-domain}
\end{equation}
The elementary turning estimate
\begin{equation}
 G_1(s)=\frac1\pi\int_s^1\arccos u\,du
 \geq\frac{2\sqrt2}{3\pi}(1-s)^{3/2}
 \label{rs:eq:G-turning}
\end{equation}
is immediate: if $v=\arccos u$, then
$1-u=1-\cos v\leq v^2/2$, so
$\arccos u\geq\sqrt{2(1-u)}$; integration over $[s,1]$ gives the displayed
constant.
It implies
\[
 \tau\geq\frac{2\sqrt2}{3}\sqrt{1-\rho}
 \geq\frac{2\sqrt2}{3}\sqrt{\frac{11}{50}}>\frac14.
\]
The last inequality follows after squaring from
$44/225-1/16=479/3600>0$.

On $(0,\tau)$, the measure $-dg$ is positive.  Hence
\begin{align}
 -\int_{(0,\tau)}\mathfrak W\,dg
 &\geq -\int_{(0,\tau)}L_\sharp\,dg \notag\\
 &=-L_\sharp(\tau)g(\tau)
   +\int_0^\tau L_\sharp'(t)g(t)\,dt.
 \label{rs:eq:R-sharp-parts}
\end{align}
There is no boundary contribution at zero.  On the outer branch
$t=A(z)=G_K(z)$ and $g(t)dt=-2z\,dz$, so
\begin{equation}
 \int_0^\tau L_\sharp'(t)g(t)\,dt
 =\int_{\rho K}^K2z\min\{G_K(z),1/4\}\,dz.
 \label{rs:eq:R-sharp-action}
\end{equation}
Since $\tau>1/4$, the interface factor becomes
\[
 g(\tau)\left(\frac\tau4+\frac3{32}-L_\sharp(\tau)\right)
 =\frac{g(\tau)}8.
\]
Lemma~\ref{rs:lem:retained} and \eqref{rs:eq:g-data} now give
\begin{equation}
 \mathcal D_{\rm rad}\geq\mathcal H_\sharp(\rho,K),
 \label{rs:eq:H-sharp-bound}
\end{equation}
where
\begin{equation}
 \mathcal H_\sharp(\rho,K):=
 \frac{\rho^2K^2}{4}
 -\frac{\pi\rho K}{4(1-\rho)}
 +\frac{\pi\rho K}{4\arccos\rho}
 +\int_{\rho K}^K2z\min\{G_K(z),1/4\}\,dz.
 \label{rs:eq:H-sharp}
\end{equation}
The simplified interface payment is asserted only on
\eqref{rs:eq:global-domain}.  Below $\tau=1/4$, its unsimplified factor is
$g(\tau)(\tau/4+3/32-\tau^2/2)$.

\subsubsection*{The universal ball quarter-layer}

Let $R_B$ be the decreasing inverse of the restriction of $G_K$ to
$[0,K]$.  Since
$G_K(\rho K)=\tau>1/4$, truncated layer cake gives
\begin{equation}
 \frac{\rho^2K^2}{4}
 +\int_{\rho K}^K2z\min\{G_K(z),1/4\}\,dz
 =B_0(K):=\int_0^{1/4}R_B(t)^2\,dt.
 \label{rs:eq:ball-layer}
\end{equation}
Therefore
\begin{equation}
 \mathcal H_\sharp(\rho,K)
 =B_0(K)-\frac{\pi\rho K}{4(1-\rho)}
 +\frac{\pi\rho K}{4\arccos\rho}.
 \label{rs:eq:H-ball}
\end{equation}

From \eqref{rs:eq:G-turning} and $G_K(z)=KG_1(z/K)$,
\begin{equation}
 R_B(t)\geq K-C K^{1/3}t^{2/3},
 \qquad
 C:=\left(\frac{3\pi}{2\sqrt2}\right)^{2/3}
   =\frac{(3\pi)^{2/3}}2.
 \label{rs:eq:ball-inverse}
\end{equation}
The right side is nonnegative for $0\leq t\leq1/4$: it is enough that
$K\geq C^{3/2}/4=3\pi/(8\sqrt2)$, whereas
$K\geq\pi/(1-\rho)>\pi$.  Squaring and integrating gives
\begin{equation}
 B_0(K)\geq\frac{K^2}{4}-\alpha K^{4/3}+\beta K^{2/3},
 \label{rs:eq:B0-lower}
\end{equation}
where
\begin{equation}
 \alpha:=\frac{6C}{5\,4^{5/3}},
 \qquad
 \beta:=\frac{3C^2}{7\,4^{7/3}}.
 \label{rs:eq:alpha-beta}
\end{equation}

Combining \eqref{rs:eq:master-defect},
\eqref{rs:eq:angular-sharp}, \eqref{rs:eq:H-sharp-bound}, and
\eqref{rs:eq:B0-lower}, and using $T\geq1$, hence $N\geq1$, gives
\begin{equation}
 \begin{aligned}
 W-P>{}&\frac{1+2\rho^2}{12}K^2
 -\alpha K^{4/3}+\beta K^{2/3}
 -\frac{\pi\rho K}{4(1-\rho)}\\
 &+\frac{\pi\rho K}{4\arccos\rho}+d(\rho).
 \end{aligned}
 \label{rs:eq:pre-scalar-gap}
\end{equation}

\subsubsection*{Exact scalar closure}

\begin{lemma}[Elementary angle bound]
\label{rs:lem:angle}
For $0\leq\rho<1$,
\begin{equation}
 \arccos\rho\leq\frac\pi2\sqrt{1-\rho}.
 \label{rs:eq:angle-bound}
\end{equation}
\end{lemma}

\begin{proof}
Write $\rho=\cos\theta$, $0\leq\theta\leq\pi/2$, and
$y=\theta/2$.  Since $\sin y/y$ decreases on $[0,\pi/4]$,
\[
 \frac{\sin y}{y}\geq\frac{2\sqrt2}{\pi}.
\]
Indeed,
$(\sin y/y)'=(y\cos y-\sin y)/y^2<0$, because
$(\sin y-y\cos y)'=y\sin y>0$.
Using $1-\rho=2\sin^2y$ gives \eqref{rs:eq:angle-bound}.
\end{proof}

\begin{lemma}[Rational bounds for $\pi$]
\label{rs:lem:rational-pi-bounds}
One has
\[
 \frac{157}{50}<\frac{333}{106}<\pi<\frac{22}{7}.
\]
\end{lemma}

\begin{proof}
Put $a=\arctan(1/5)$ and $b=\arctan(1/239)$.  The tangent addition
formula gives
\[
 \tan(2a)=\frac5{12},\qquad
 \tan(4a)=\frac{120}{119},\qquad
 \tan(4a-b)=1.
\]
Here $b<a$, while $(6/5)^2<2$ gives
$1/5<\sqrt2-1=\tan(\pi/8)$, and hence $a<\pi/8$.  Therefore
$0<4a-b<\pi/2$.  Thus this is Machin's identity
$\pi=16a-4b$.  The alternating arctangent expansion gives
\[
 a>\frac15-\frac1{3\,5^3}+\frac1{5\,5^5}-\frac1{7\,5^7},
 \qquad b<\frac1{239}.
\]
Indeed, these inequalities follow by integrating
\[
 \frac1{1+t^2}=1-t^2+t^4-t^6+\frac{t^8}{1+t^2}
 \qquad(0\leq t\leq1)
\]
and using $1/(1+t^2)<1$ for $t>0$.
Thus
\[
 \pi>
 \frac{1231847548}{392109375}
 =\frac{333}{106}+\frac{3418213}{41563593750},
 \qquad
 \frac{333}{106}-\frac{157}{50}=\frac2{1325}>0.
\]
For the upper bound,
\[
 \frac{x^4(1-x)^4}{1+x^2}
 =x^6-4x^5+5x^4-4x^2+4-\frac4{1+x^2},
\]
so integration and $4\arctan1=\pi$ give
\[
 0<\int_0^1\frac{x^4(1-x)^4}{1+x^2}\,dx
 =\frac{22}{7}-\pi.
\]
Also \(22/7-333/106=1/742>0\), completing the asserted chain.
\end{proof}

Put $e:=1-\rho$.  It suffices to prove positivity of
\begin{equation}
 \begin{aligned}
 \underline{\mathcal G}(\rho,K):={}&
 \frac{1+2\rho^2}{12}K^2
 -\alpha K^{4/3}+\beta K^{2/3}
 -\frac{\pi\rho K}{4e}\\
 &+\frac{\rho K}{2\sqrt e}
 +\frac3{4\sqrt e}-\frac14.
 \end{aligned}
 \label{rs:eq:G-under}
\end{equation}
Define
\begin{equation}
 \mathsf A:=\alpha\pi^{4/3},
 \quad \mathsf B:=\beta\pi^{2/3},
 \quad C_1:=\frac43\alpha\pi^{1/3},
 \quad C_2:=\frac23\beta\pi^{-1/3}.
 \label{rs:eq:scaled-constants}
\end{equation}
Positive-side cubing and Lemma~\ref{rs:lem:rational-pi-bounds} give
\begin{equation}
 \mathsf A<\frac{49}{40},
 \qquad
 \frac{179}{1000}<\mathsf B<\frac15,
 \qquad
 C_1<\frac{13}{25},
 \qquad
 C_2<\frac1{25}.
 \label{rs:eq:constant-bounds}
\end{equation}
For completeness,
\begin{equation}
 \mathsf A^3=\frac{3^5\pi^6}{5^3 2^{10}},
 \quad
 \mathsf B^3=\frac{3^7\pi^6}{7^3 2^{20}},
 \quad
 C_1^3=\frac{3^2\pi^3}{5^3 2^4},
 \quad
 C_2^3=\frac{3^4\pi^3}{7^3 2^{17}}.
 \label{rs:eq:constant-cubes}
\end{equation}
Thus each comparison in \eqref{rs:eq:constant-bounds} is an exact rational
cross multiplication after cubing.

\paragraph{Exact scalar reserve table.}
For reproducibility, the positive margins used in those cubed comparisons
and in the final convexity propagation are recorded below.  The first five
rows are respectively the rational difference between the cube of the stated
upper bound and the certified upper cube, or between the certified lower cube
and the cube of the stated lower bound.
\begin{center}
\small
\begin{tabular}{@{}ll@{}}
\toprule
inequality&exact positive reserve\\
\midrule
\(\mathsf A<49/40\)&\(13125773/1505907200\)\\
\(\mathsf B>179/1000\)&
\(521715693010963/5619712000000000000\)\\
\(\mathsf B<1/5\)&\(176853008713/82644187136000\)\\
\(C_1<13/25\)&\(9767/10718750\)\\
\(C_2<1/25\)&\(243014341/30118144000000\)\\
turning-interface square above \(1/16\)&\(479/3600\)\\
\(F_0'''\) lower bound&\(676141/2450\)\\
\(D_0''\) lower bound&\(1951/100\)\\
\(\partial_K^2\underline{\mathcal G}\) lower bound&
\(47447/443682\)\\
\bottomrule
\end{tabular}
\end{center}

Set
\begin{equation}
 K_0:=\frac\pi e,
 \qquad x:=e^{-1/6}\geq1.
 \label{rs:eq:base-x}
\end{equation}
Direct substitution gives
\begin{equation}
 \underline{\mathcal G}(\rho,K_0)=F_0(x),
 \label{rs:eq:G-base}
\end{equation}
where
\begin{equation}
 \begin{aligned}
 F_0(x):={}&\frac\pi2x^9-\mathsf A x^8
 -\frac{\pi^2}{12}x^6+\mathsf Bx^4\\
 &+\left(\frac34-\frac\pi2\right)x^3
 +\frac{\pi^2}{6}-\frac14.
 \end{aligned}
 \label{rs:eq:F0}
\end{equation}
The constant bounds imply
\begin{equation}
 F_0(1)>\frac{8269}{30000}>\frac14,
 \quad
 F_0'(1)>-\frac{14437}{6125}>-\frac52,
 \quad
 F_0''(1)>\frac{18162}{1225}>14.
 \label{rs:eq:F0-endpoint}
\end{equation}
For $x\geq1$, group the leading terms as
$x^5(252\pi x-336\mathsf A)$ and use $x^5\geq x^3$ after the
resulting coefficient is positive; indeed its rational lower bound is
\[
 252\frac{157}{50}-336\frac{49}{40}=\frac{9492}{25}>0.
\]
Dropping the positive
$24\mathsf Bx$ term then gives
\begin{equation}
 \begin{aligned}
 F_0'''(x)
 &>x^3\left(
 252\frac{157}{50}-336\frac{49}{40}
 -10\left(\frac{22}{7}\right)^2\right)-\frac{69}{14}\\
 &\geq\frac{676141}{2450}>0.
 \end{aligned}
 \label{rs:eq:F0-third}
\end{equation}
Hence $F_0''(x)>14$.  With $y=x-1\geq0$, Taylor's formula yields
\begin{equation}
 F_0(x)>\frac14-\frac52y+7y^2
 =7\left(y-\frac5{28}\right)^2+\frac3{112}>0.
 \label{rs:eq:F0-positive}
\end{equation}

A direct differentiation gives
\begin{equation}
 \partial_K\underline{\mathcal G}(\rho,K_0)=D_0(x),
 \label{rs:eq:base-derivative}
\end{equation}
where
\begin{equation}
 D_0(x):=
 \frac\pi{12}(3x^6-5+4x^{-6})
 +\frac12(x^3-x^{-3})-C_1x^2+C_2x^{-2}.
 \label{rs:eq:D0}
\end{equation}
At $x=1$,
\begin{equation}
 D_0(1)>\frac1{300},
 \qquad D_0'(1)>\frac{54}{175}.
 \label{rs:eq:D0-endpoint}
\end{equation}
Moreover,
\begin{equation}
 \begin{aligned}
 D_0''(x)={}&\frac\pi{12}(90x^4+168x^{-8})
 +3x-6x^{-5}-2C_1+6C_2x^{-4}\\
 &>\frac{15}{2}\frac{157}{50}-3-\frac{26}{25}
 =\frac{1951}{100}>0.
 \end{aligned}
 \label{rs:eq:D0-convex}
\end{equation}
Here $x\geq1$ gives $90x^4\geq90$ and
$3x-6x^{-5}\geq-3$; the omitted terms are nonnegative, and
$-2C_1>-26/25$.  Thus the displayed lower bound follows term by term.
Since $D_0''(x)>0$,
\[
 D_0'(x)\geq D_0'(1)>0\qquad(x\geq1).
\]
Thus $D_0$ is increasing and $D_0(x)\geq D_0(1)>0$.

Finally,
\begin{equation}
 \partial_K^2\underline{\mathcal G}
 =\frac{1+2\rho^2}{6}
 -\frac{4\alpha}{9}K^{-2/3}
 -\frac{2\beta}{9}K^{-4/3}.
 \label{rs:eq:G-K-convex}
\end{equation}
For $K\geq K_0\geq\pi$,
\[
 \alpha K^{-2/3}\leq\frac{\mathsf A}{\pi^2},
 \qquad
 \beta K^{-4/3}\leq\frac{\mathsf B}{\pi^2}.
\]
Consequently,
\begin{equation}
 \partial_K^2\underline{\mathcal G}
 >\frac16-
 \frac{4(49/40)+2(1/5)}{9(157/50)^2}
 =\frac{47447}{443682}>0.
 \label{rs:eq:G-K-positive}
\end{equation}
Since \eqref{rs:eq:G-K-positive} is positive,
\[
 \partial_K\underline{\mathcal G}(\rho,K)
 \geq\partial_K\underline{\mathcal G}(\rho,K_0)
 =D_0(x)>0
 \qquad(K\geq K_0).
\]
Hence
\begin{equation}
 \underline{\mathcal G}(\rho,K)
 \geq\underline{\mathcal G}(\rho,K_0)
 =F_0(x)>0
 \qquad(K\geq K_0).
 \label{rs:eq:G-positive}
\end{equation}

\begin{theorem}[Global low- and middle-ratio proxy theorem]
\label{rs:thm:global-proxy}
For
\[
 0<\rho\leq\frac{39}{50},
 \qquad K\geq\frac\pi{1-\rho},
\]
one has
\begin{equation}
 P(\rho,K)<W(\rho,K).
 \label{rs:eq:proxy-conclusion}
\end{equation}
Consequently,
$N_D(A_{\rho,1},K^2)<W(\rho,K)$.
\end{theorem}

\begin{proof}
Equation \eqref{rs:eq:pre-scalar-gap}, the angle bound, and
\eqref{rs:eq:G-positive} give
\[
 W-P>\underline{\mathcal G}(\rho,K)>0.
\]
The layer-cake and angular-block proofs retain every strict wall.
\end{proof}
% END mechanically inlined support source

\end{document}
